\scp{Wallacea and its peopling}{02-01}
Before focusing on the prehistory
and classification of the Austronesian languages of the region,
it is helpful to understand the context.
There are two biogeographical regions with relatively shallow seas connecting islands within each of
them, and relatively deep water crossings between the two zones.
\it{Sundaland} extends from mainland South\ili{east} Asia as far \ili{east} as the islands of Java, Bali, and Borneo.
\it{Sahulland} includes Australia, New Guinea,
and the \ili{Aru} Islands in southeastern Maluku.
(\frf{02-01}.)

\begin{figure}
	\includegraphics[width=\columnwidth]{WallaceLine.pdf}
	\fcp{Biogeographical zones with Wallacea as a transition zone}{02-01}
\end{figure}

The \it{Wallace Line} in \frf{02-01} indicates the westernmost extent of fauna (such as marsupials) commonly associated with Sahulland.
The \it{Lydekker Line} runs along the border of the Sahul continental shelf.
\it{Weber’s Line} is variously described as the area where the two types of fauna (from Sundaland
and Sahulland) are thought to be equally mixed,
or as the easternmost extent of certain categories of fauna from Sundaland (such as reptiles).
The flora (plants) do not show such clear patterns.
Wallacea is the transition zone between the two biogeographical regions (the zone between the Wallace
Line and the Lydekker Line),
and is a zone that requires multiple deep water crossings for animals or humans to traverse
\citep[58]{Bellwood2007,Bellwood2019ed,OConnor2015,Oppenheimer2003}.
We return to this in \crf{ch:MorLes}.

Wallacea includes flora and fauna from both Sundaland
and Sahulland, as well as flora
and fauna endemic (native to,
unique) in the region such as the \it{babirusa} (\it{Babyrousa babyrussa;} \frf{02-02})
and the \ili{Buru} red lory (\it{Eos bornea cyanonothus} \frf{02-RedLory}).
Wallacea is where Austronesian-speaking peoples coming from the north or west would have first encountered
marsupials, and developed a need to talk about them \citep{Blust1982,Blust2012b,Schapper2011}.
Biogeographical Wallacea includes Sulawesi,
Lombok, and western \ili{Sumbawa}.
“Linguistic Wallacea” does not include Sulawesi,
Lombok, or western \ili{Sumbawa}.
Linguistic Wallacea excludes these western parts
and extends further eastward to encompass parts of Indonesian Papua, including SHWNG languages.
Our study is focused on linguistic Wallacea minus SHWNG
(see Figures \ref{fig:01-05}, \ref{fig:01-06}, \ref{fig:02-01}).

%\begin{figure}
	%\includegraphics[width=0.4\columnwidth]{02-02.jpg}
	%\fcp{Babirusa (\it{Babyrousa babyrussa})}{02-02}
%\end{figure}

\begin{figure}\centering
  \subfigure[Babirusa (\it{Babyrousa babyrussa}) ]{\label{fig:02-02}
    \includegraphics[width=0.45\textwidth]{02-02.jpg}
    }\hspace{15mm}
    \subfigure[\ili{Buru} red lory (\it{Eos bornea cyanonothus})]{\label{fig:02-RedLory}
    %\subfigure[\ili{Buru} red lory]{\label{fig:02-RedLory}
    \frame{\includegraphics[width=0.25\textwidth]{02-RedLory.jpg}}
    }
  \caption{Fauna unique to Wallacea}
\end{figure}

Humans travelling from Asia to reach Australia
and New Guinea, regardless of whether they came from the north or west,
would have had to make multiple deep water crossings in Wallacea.
\citet[1]{Bellwood2019b} points out,
“No land bridges ever crossed the full expanse of this region of deep seas
and steeply plunging coastlines during the Pleistocene Ice Ages,
even during glacial maxima,
and humans migrating from Taiwan or Sundaland towards New Guinea
and Australia always had to cross sea gaps between islands to reach their goals.” 

At the time of initial human encroachment into the region,
\citet[215]{OConnor2015} is even more specific,
“the global sea level was 60--70 m lower than it is today.
While this may have reduced distances between islands,
at l\ili{east} eight separate crossings,
and one leg of at l\ili{east} 70 km,
would have had to be made,
regardless of the route taken” to reach Australia.
(See \frf{02-04}.)

\begin{figure}
	\includegraphics[width=\columnwidth]{02-04.jpg}
	\fcp{Deep water crossings in Wallacea}{02-04}
\end{figure}

\citet{BellwoodIrwin1998} find evidence of human habitation in north Maluku around 35,000 BP (BP = before present).\footnote{
		Technically BP is set at 1950,
		when the term “before present” came into use.}
At greater depths,
around the 40,000 BP range,
evidence of human habitation has not been found there \citep{Bellwood2019b}.

\citet[217]{OConnor2015} reports that in modern day Timor-Leste,
“at two sites, Jerimalai shelter and Lene Hara Cave,
[human] occupation dates back to 42,000 cal BP”. %(cal = calibrated).
She also reports:

\begin{quote}
	``More than thirty-nine thousand fish bones from twenty-three
	different taxa were identified overall from this test pit.
	The most common fish were Scombrids (tuna),
	Scaridae (parrotfish), Carangidae (trevallies),
	Balistidae (triggerfish), and groupers (Serranidae).
	Other major fish families identified are Lethrinidae (emperors),
	Lutjanidae (snappers), unicornfish (Acanthuridae),
	wrasses (Labridae), and puffers (Tetradontidae).
	Rays and sharks (Elasmobranchi) are also present \citep{OConnorClarkson2011}.
	Next to fish, marine turtles made up a large part of the diet.
	[{\ldots}] Although fish remains are most plentiful in Holocene levels,
	the earliest Pleistocene occupation level (42,000
	and 38,000 years ago) contained significant quantities of fish bone.
	[{\ldots}] Fast-swimming pelagic species [deep sea,
	open ocean], such as tunas, account for almost 50 percent
	of the fish bones in these levels.''
	\citep[219]{OConnor2015}
\end{quote}

She notes that to reach Australia or New Guinea from Sundaland,
“Whatever route was taken,
this was a remarkable accomplishment at this early date
and demonstrates the capacity for complex planning
and technological innovation that is the hallmark of our species.
First, it required the successful navigation of multiple water crossings between the stepping-stone islands en route.
Second, the colonists arriving on each new island had to adapt to unfamiliar fauna,
flora, and landscapes” \citep[215]{OConnor2015}.

\begin{figure}\centering
	\includegraphics[width=0.4\columnwidth]{02-03.jpg}
	\fcp{Skipjack tuna (\it{Katsuwonus pelamis})}{02-03}
\end{figure}

The implications of the large numbers of pelagic
fish bones (see \frf{02-03}) are highly significant for
our understanding of the prehistory of the region.

\begin{quote}
	``The pelagic fish bones from Jerimalai clearly demonstrate that the level of maritime capacity
	and fishing technology of the early modern humans
	who colonised our region was much higher than previously believed.
	[{\ldots}] Pelagic species are fast-moving,
	travel long distances, and are mostly caught offshore.
	They are therefore quite difficult to capture,
	as compared with inshore fish.
	[{\ldots}] Whatever method was used,
	the high proportion of pelagic species at Jerimalai indicates
	that their capture was probably planned rather than fortuitous.
	{\ldots} The earliest direct evidence of fishing technology in the East Timor sites is the tip of a hook
	made on Trochus sp. shell found in Jerimalai in a level dated to between 23,000
	and 16,000 years ago \citep{OConnorClarkson2011}.
	Indeed, this currently stands as the earliest evidence for the manufacture of fishhooks anywhere
	in the world.'' \citep[220f]{OConnor2015}
\end{quote}

She concludes:

\begin{quote}
	``The early modern humans moving out of mainland Asia must already have possessed a high degree of maritime
	capacity, which would have enabled them to make water crossings into Island South\ili{east} Asia in the first place.
	Settlement in the islands of Wallacea would have presented many new challenges,
	not l\ili{east} of which would have been the restricted range of terrestrial animal species on many of them.
	[{\ldots}] by 42,000 cal BP,
	the people living in this region were skilled in the acquisition of marine resources,
	which requires high levels of planning.
	They caught fast-moving pelagic fish such as tuna,
	as well as sharks,
	rays, and a range of inshore fish; to do this,
	they must have possessed a range of sophisticated organic technology.''
	\citep[221]{OConnor2015}
\end{quote}

We do not know what languages these earliest human travellers into
and through the region spoke.
We do know that they were not Austronesian languages at that early time.
We do know that early humans did reach
and settle the New Guinea region far to the \ili{east},
requiring many deep water crossings.
We do not know if they had a continuous presence throughout our target region or only in some places.
For example, \citet[220]{OConnor2015} says, “Jerimalai shelter appears to have little evidence of occupation
38,000--24,000 years ago {\ldots} [but] nearby Lene Hara Cave was occupied during this period.”\footnote{
		The relative gaps in dating are interesting,
		and probably reflect the scarcity of archeological digs in the Wallacea region.
		Presumably early humans travelling to Sahulland (the Australia
		and New Guinea region) would have had to traverse Wallacea.
		Dating human habitation in Sahulland currently ranges
		from 50,000--60,000 BP \citep{OConnellAllen2003,ReddStoneking1999}.
		Yet the earliest dates in Wallacea seem to be around 42,000 BP.
		One assumes that additional digs will find evidence of earlier dating,
		or possibly evidence of other routes to Sahulland.}

By the time the Austronesians arrived in the region
tens of thousands of years later, around 2,000--1,500
BC (\citeauthor{Bellwood2007} [\citeyear[123]{Bellwood2007}]; ≈ 4,000--3,500 BP),
it appears that the pre-existing populations already in place spoke various non-Austronesian languages
of at l\ili{east} two different language families that are loosely combined under the term ‘Papuan’.

Before returning to discussion of the pre-existing non-Austronesian languages of eastern
Indonesia and Timor-Leste, it is helpful to draw attention to a few points
regarding evidence of the arrival of Austronesian-speaking peoples.

\begin{itemize}
	\item There is abundant evidence of the introduction of various flora (particularly
				food crops) and fauna associated with the arrival of the Austronesian-speaking peoples
				into the region around 4,000-3,500 BP \citep{Bellwood2007}.
	\item There is evidence of “Neolithic settlement of northern Maluku”,
				associated with the Austronesian-speaking peoples dating about 3,500 BP \citep[514]{Spriggs2011}.
				More specifically “plain red-slipped pottery,
				polished stone adzes, pigs,
				and dogs dating from 3300 BP in the Uattamdi rockshelter”
				assumed to be associated with Austronesian-speaking
				peoples have been found on Kayoa Island \citep[218]{Bellwood2019a}.
	\item There are traces of “pig-bone [{\ldots}] red-slipped pottery,
				1 sherd with incised decoration; chert
				and obsidian flakes” associated with Austronesian-speaking peoples at a site in the \ili{Banda} Islands in
				central Maluku dating around 3,000 BP \citep[520--521]{Spriggs2011}.
	\item It is entirely possible that the Austronesian incursion(s) into northern Maluku are separate from
				and independent of the incursions into the Lesser Sundas,
				southern and central Maluku (see \srf{15-06}).
	\item There was a period of intense,
				rapid, and far-reaching migrations by the Austronesian-speaking peoples,
				covering thousands of islands
				and thousands of kilometres spanning just a few centuries.
				\citet[520--521]{Spriggs2011} says,
				“The spread of Neolithic \tsc{An}-speaking cultures across much of ISEA
				[Island South\ili{east} Asia] is a similar phenomenon,
				in terms of its rapidity,
				to the Lapita expansion beyond the Bismarck Archipelago between about 3100
				and 2900 BP when that culture spread beyond Near Oceania
				through the south-\ili{east} Solomons,
				Vanuatu, New Caledonia, Fiji
				and into western Polynesia.
				[{\ldots}] Polynesia, which was first occupied only
				about 3500 BP in the case of the Marianas
				and possibly Palau, and not long after about 3100 BP
				for the rest of Island Melanesia
				and western Polynesia.”
\end{itemize}

So, we note that Remote Oceania appears to have been settled by people speaking Austronesian languages
around a broadly similar time period as for the settlement of eastern Indonesia.
It is not clear if this rapid expansion into the Pacific involved or by-passed what is now central
and southern Maluku, Timor-Leste,
and the rest of Nusa Tenggara (the Lesser Sundas).
The assumption based on the distribution of EMP languages (including SHWNG) has been that the jump-off
point for what later became Oceanic
and Polynesian languages somehow involved north Maluku.
But that view may no longer have strong support from archaeology \citep[4]{Bellwood2019b}.
We return to this issue in \crf{ch:MorLes}.

We do know that many grammatical features commonly associated with Oceanic languages
\citep{Pawley1973,Ross1998} are also found in many languages
of present day eastern Indonesia and Timor-Leste.
These include the grammar and semantics of active
and non-active verbs, verbs of motion
and posture, serial verbs,
prepositional verbs, complex possessive constructions, alienable (indirect)
and inalienable (direct) possession.
And, for example, an edible possessive construction commonly associated with some Oceanic languages,
is also found in \ili{Selaru} in southern Maluku \citep[47--48]{Coward2005}.
Phonological features associated with Oceanic languages such as PMP *-ay > \it{e},
*-aw > \it{o}, and *ə > \it{o} are found sporadically in different parts of the target region.
But these are also found widely in languages in Sulawesi
and other languages west of the Blust line.

In trying to see where the linguistic picture
and the archaeological picture link up,
\citet[511]{Spriggs2011} comments: 

\begin{quote}
	``We can use the spread of the ISEA Neolithic as a proxy for AN language spread,
	as justified at length by \citet{Pawley2004}
	and \citet{Ross2008}, among others.
	\it{In doing this, it is very hard to see anything between PMP
	and EMP at all from the archaeology.} It would seem that movements out of Taiwan were rapid after about
	4000 BP and by 3800 BP dialects of PMP were spoken everywhere from the Philippines to eastern Borneo,
	Sulawesi and south to East Timor,
	spreading with the first pottery-using cultures in those areas.
	Currently the dates for the EMP area in northern Maluku do seem to reflect a later time of spread,
	at about 3500 BP as with Palau
	and the Marianas and Java.'' \citep[511]{Spriggs2011}	[emphasis ours]
\end{quote}

So, in terms of the limited archaeological data currently available,
Spriggs sees that the Austronesian incursion into the
Timor region occurred several hundred years before
the evidence of Austronesians in northern Maluku,
precisely during a time period where a span of
200--300 years appears to have been incredibly significant
for the dispersal of Austronesian-speaking peoples.\footnote{
		Given the scarcity of archaeological digs
		in the region, this may be making much ado about nothing.} 

There are questions about where
and when and how Austronesian-speaking peoples came
into the target region that are introduced briefly
here, and addressed in greater detail in \crf{ch:MorLes}.
\citet[98]{Blust2008} says,

\begin{quote}
		``\ili{As} noted in \citet{Blust1993},
		\ili{Central} Malayo-Polynesian almost certainly was the product of rapid movement from north\ili{east} Indonesia
		through the central and southern Moluccas and into the Lesser Sunda islands,
		forming a dialect chain that favored the ready diffusion of innovations across incipient language boundaries,
		and hence destroying any tree-like structure that may have previously arisen.'' \citep[98]{Blust2008}
\end{quote}

While he provides no evidence for such movement,
he nevertheless criticises \citet{Schapper2011} for an “implicit claim” that both the glosses in the
languages and the limited distribution of bandicoots show the Austronesian-speaking peoples encountered the marsupial
cuscus before they encountered the bandicoot,
saying her scenario “implies an entry via Sulawesi into the Lesser Sundas,
with subsequent eastward movement into the Moluccas” \citep[274]{Blust2012b}.
While that goes beyond what she actually says,
it is nevertheless revealing.
Since Blust rejects an incursion via Sulawesi,\footnote{
		We show in \crf{ch:Tal} that there is evidence
		from the historical phonology for at l\ili{east} one incursion into Maluku from Sulawesi.
		\citet[130]{Bloyd2020} also observes,
		“Sula was a casual swim’s distance from Sulawesi during the height of the Ice Age.”
		In chapters \ref{ch:Mar}
		and \ref{ch:MorLes} we also show several lexical ties between south\ili{east} Sulawesi,
		Maluku, and the rest of Wallacea.}
		we have to assume by “north\ili{east} Indonesia” (\citeyear[98]{Blust2008})
		he means Halmahera or Indonesian Papua.
Yet the limited archaeological evidence currently available has the Austronesian-speaking peoples reaching
Timor (presumably from the west),
300-500 years before they show up in Halmahera.
The scattered and disjointed pockets of conservative features in the historical phonology,
along with common direction of sound changes presents a much more complex picture than that proposed by Blust.

The scattered and uneven pockets of complex morphology along with non-cognate morphology from subgroup
to subgroup that does not trace back to PMP,
interspersed with other pockets of fairly simple morphology,
also paints a much more complex picture than that proposed by Blust.
And in a context of great mobility along with intense
and prolonged contact with typologically different languages,
the principle of “the area with the greatest diversity probably represents the homeland” cannot be assumed to hold true.
So while Blust’s scenario of Austronesian incursion into our target region does not seem to be grounded
in archaeology, historical phonology, morphosyntax, or the lexicon, we find that the linguistic evidence favours multiple
incursions, perhaps from different starting points,
and does not necessarily rule out his proposed scenario as one of several possibilities.
Yet we are also not actually sure what data or evidence his proposed scenario is based on.
We favour a both-and approach over an either-or approach on Austronesian incursions into the region.
We return to this issue in \crf{ch:MorLes}.

\citet{Spriggs2011} likewise comments on the broad patterns of contact by the Austronesian-speaking peoples
as they encountered other populations,
and implications for methodology.

\begin{quote}
	``Roger Green’s \citeyear{Green1991,Green2000} model of intrusion,
	integration and innovation captures the situation well in ISEA as well as the western Pacific,
	whether we are talking of material culture, language or people.
	There was indeed some migration out of Taiwan \citep{Kayser2008};
	there was mass recruitment of people from populations already resident in ISEA
	and Near Oceania as the Neolithic spread \citep{Soares2011};
	artefacts and practices were integrated from already-resident groups
	and others were discarded by them; new ideas were brought into being as unexpected human
	and environmental situations were encountered.'' \citep[523]{Spriggs2011}

	``[The] Neolithisation of ISEA was a new process of identity formation that seized the imagination of a
	mass of people on hundreds of islands across thousands of kilometres of ocean,
	spreading like a pulse across ISEA
	and into the Pacific over a few centuries.
	It spread through processes both of migration and recruitment in-place.''
	 \citep[524]{Spriggs2011}

	``[Consequently:] we are struggling to find appropriate models to deal with just what happens during such
	temporal ‘hot spots’.'' \citep[523]{Spriggs2011}									
\end{quote}

Contact complicates comparative-historical linguistics.
And contact is not a trivial add-on,
but rather a core feature of our target region,
and one that is necessary to understand for unravelling
the complexities of the data in eastern Indonesia and Timor-Leste.
Many people are finding the traditional branching tree model (see \frf{01-02})
and an assumption of a monolithic spread of the
Austronesian-speaking peoples to be inadequate to account
for the complexities of the region.
\citet[9]{Klamer2019} summarises the issues:

\begin{quote}
	``Indeed, for the micro-groups that have been studied reasonably well (cf.
	\citealp{Adelaar2005}), the absence of evidence for intermediate
	ancestral nodes connecting them suggests a complex dispersal scenario.
	There may have been just a few MP groups migrating into ISEA,
	but then each of these groups did not spend a sufficiently long period of time together as unified speech
	communities after their split from MP.
	Because if they had,
	they would have innovated words in their lexicon or developed different pronunciations which would have
	been inherited in their daughter languages
	and could now be used to define the subgroup as such by its unique lexical and phonetic innovations.
	There may also have been not a few but \it{many} different MP communities migrating at different times
	and in different directions throughout history.
	Overall, the lack of shared inherited features suggests that most of the MP languages that dispersed
	into ISEA shared a relatively short local micro-history.
	In contrast, the shared inherited features of,
	for example, proto-Oceanic, indicate that this group did remain a unified speech community for a period
	that was sufficiently long to create such defining innovations
	and pass them on to their daughter languages.
	In sum, there is no evidence that the dispersal of MP speakers in ISEA followed simple trajectories;
	instead, it is more likely that multidirectional processes of contact
	and language admixture have been at work for millennia.''
	\citep[9]{Klamer2019} [emphasis in the original]
\end{quote}

Returning now to the Papuan languages,
we note that those still spoken in the region today include the \tsc{Timor-\ili{Alor}-Pantar} (TAP) languages
(\citealp{Klamer2017,HoltonRobinson2017} -- see \frf{02-05}),
and the \tsc{West Papuan} languages of North Halmahera
(\citealp{Voorhoeve1988,GrimesGrimes1994,HoltonKlamer2018} -- see \frf{02-06}).
But the implications of \citet{Donohue2007} are also that we have evidence of earlier Papuan languages
spoken as far west as \ili{Sumbawa} into the early 1800s.
In this case, the \ili{Tambora} language was wiped out,
not by the Austronesian-speaking peoples,
but by an erupting volcano.
There is widespread trace evidence of various sorts of contact with SOV (Papuan) languages found
in Austronesian languages that are hundreds of kilometres from any Papuan languages spoken today,
such as on Flores,
Rote, \ili{Buru}, Seram, \ili{Kei}, Tanimbar, and \ili{Babar}.
We elaborate on these issues in \srf{02-03-01}.

\begin{figure}[ht]
	\includegraphics[width=\columnwidth]{TimorPapuan.pdf}
	\fcp{Austronesian and Papuan languages in greater Timor}{02-05}
\end{figure}

\begin{figure}[p]
	\centering
	\includegraphics[width=0.87\columnwidth]{HalmaheraPapuan.pdf}
	\fcp{Austronesian and Papuan languages in north Maluku}{02-06}
\end{figure}

There are interesting questions about the prehistory
of the non-Austronesian languages in the region.
(\frf{02-05} and \frf{02-06}.)
There are several obvious possibilities:

\it{Scenario {\#}1:} The non-Austronesian languages are remnants of the ancestral languages spoken by
the earliest humans on the way from Island South\ili{east} Asia to the island of New Guinea.
These ancestral languages were perhaps the linguistic precursors of the many Papuan language families
that later diversified significantly.

\it{Scenario {\#}2:} The non-Austronesian languages today are the result of more recent westward moving
back migration(s) out of the New Guinea mainland.
The assumption in this scenario,
based partly on the spread from the New Guinea region of several important food crops
and animals, is that they were already in place in the Timor region before the arrival of the Austronesian-speaking
peoples around 4,000--3,500 BP.
This view is reflected in \citet[85]{Blust2013} with regards to the TAP languages,
“All non-AN languages of
the Lesser Sundas have been assigned to the Trans-New Guinea Phylum of Papuan languages,
and evidently represent a prehistoric back-migration from the western end of New Guinea to the islands
of \ili{Pantar}, \ili{Alor} and adjacent parts of Timor.”\footnote{
		In spite of accepting such a scenario of a prehistoric back-migration
		from western New Guinea,
		\citet[40]{Blust2009} appears sceptical that “the reference to
		‘contact-induced change’ draws on several publications
		by Donohue that advance the novel theory that a Papuan-based
		trading network spanned the \ili{Malay} archipelago
		prior to the arrival of AN speakers during the second millennium BC.”}
There is a growing body of typological
and lexical evidence of a connection between TAP
and \tsc{West Bomberai} languages which are thought to link to \tsc{Trans New Guinea} languages
\citep{Windschuttel2018,UsherSchapper2018,UsherSchapper2022}.

\it{Scenario {\#}3:} a combination of scenarios {\#}1 + {\#}2.
While there is increasing evidence in support of scenario {\#}2,
that nevertheless does not negate the possibility of scenario {\#}1.
That is, the more recent westward back-migration from the New Guinea region intermixed with a remnant
Papuan substrate from the early human migrants who left evidence in the archaeology.

The reality is actually much more complicated than either {\#}1 or {\#}2.
Given the span of time of over 40,000 years,
and the evidence of early boat technology that supported deep water crossings
and open ocean fishing,
it is unrealistic to limit non-Austronesian incursions to one time,
in one place and from one direction.

For example, there is evidence of at l\ili{east} two different waves of non-AN languages in \ili{Alor}-\ili{Pantar}.
Several groups claim their ancestors came from the Timor mainland.
The Papuan language enclave of \ili{Oirata} on \ili{Kisar} Island \ili{east} of Timor also traces their origins to a relatively
recent migration from the Timor mainland.
Some of these could have happened at the
same time or after Austronesian contact
(\citealp[8f]{GrimesTherikGrimesJacob1997}, \citealp{Klamer2017}).

But there are many other complexities to add into the mix as well.
Scenario {\#}1 above is problematic,
in that the break-up of the \ili{Alor}-\ili{Pantar} (AP) branch of TAP appears to have occurred about the same time
as the arrival of the Austronesian-speaking peoples.
\citet[10]{Klamer2017} says,

\begin{quote}
“the level of lexical
and grammatical similarity in the AP family does not support
an age of more than several millennia,
and the reconstructed vocabulary of proto-AP appears
to contain Austronesian loan words such as ‘betel
nut’ \citep{HoltonSchapper2012,Robinson2015}.
Ancient Austronesian loans found across the \ili{Alor}-\ili{Pantar} family
following regular sound changes suggest
that the AP family split up after being in contact
with the Austronesian languages in the area.”
\end{quote}

Presumably
the higher-level relationships of the TAP
family predated the arrival of the Austronesians.
In \srf{02-05} we discuss potential methodological
problems with linking the timing of the break-up
of the AP branch of TAP to the arrival
of the Austronesian-speaking peoples.

An issue that supports scenario {\#}2 above is growing evidence of westward movement out of the New Guinea
region of certain food crops whose botanical origins come from there,
and that the dispersal was happening prior to the arrival of the Austronesian-speaking peoples
\citep{DenhamDonohue2009,DonohueDenham2010,Schapper2017}.
For example, “the initial stages of banana
and sugarcane domestication with subsequent diffusion westward [out of New Guinea] occurred before the
appearance of Malayo-Polynesian languages in ISEA” \citep[20]{DonohueDenham2010}.
They are specifically referring here to banana cultivars developed from \it{Musa acuminata} spp.
\it{banksii}, and to \it{Saccharum officinarum} (sugarcane),\footnote{
		Since sugarcane (\it{Saccharum
		officinarum}) is now known to have originated from the New Guinea region,
		it is assumed that the \tsc{PAn} reconstruction of *CəbuS ‘sugarcane’ either reflects a very early spread north
		to Taiwan, or a different \it{Saccharum} species \citep[515]{Spriggs2011}.
		Very likely this northern variety was \it{Saccharum sinense}.}
but their discussion also includes taro, breadfruit and sago.
This dispersal of food crops westward from the New Guinea region has huge implications for methodology
that we return to later.

With regard to genetic DNA research in the region,
\citet[18]{Klamer2019} summarises recent findings:

\begin{quote}
	``Genetic research confirms that during the Holocene period in ISEA,
	groups with Asian ancestry spread through ISEA.
	Although the data on ISEA are still very sparse,
	there is presently statistical support that Taiwan played
	an instrumental role in Holocene dispersals
	across ISEA, but the data do not clearly support a \it{majority}
	genetic input from there.
	[{\ldots}] \citet{HudjashovCox2017} show
	that during the period in question (<4,000 years)
	there is very limited evidence for autosomal
	gene transfer to Indonesian populations,
	and that which is tends to be from the Philippines.
	Moreover, there is no evidence of gene transfer to western Indonesia,
	and Taiwan does not seem to have a direct role.
	Overall, the speakers of Austronesian languages today
	are very diverse in biological terms,
	and the genetic evidence indicates that the evolutionary processes
	in ISEA have been very complex
	and have involved a great deal of admixture.
	Finally, the genetic data are often inconclusive about the directions
	of these dispersals'' \citep[265]{LansingSudoyo2011}.
	[emphasis in the original]
\end{quote}

\citet[22]{Fricke2019} contributes additional information
relevant to the region:

\begin{quote}
	``There appears to be a clear boundary between Asian genomes in the west
	and non-Asian genomes in the \ili{east} between the islands of Bali
	and Flores, more specifically between \ili{Sumbawa}
	and Flores 
	(\citealp[1837f]{KarafetHammer2010}, \citealp[2447]{HudjashovCox2017}).
	This genome boundary coincides with the Wallace line,
	a boundary of transition between Asian
	and Australian types of flora and fauna.
	[{\ldots}] The sampling results from the eastern part of Flores (Bama),
	from Lembata and from Timor show an almost half-half division of the Papuan
	and the Austronesian genome component.
	This genetic profile is very similar to the sampling results on the islands of \ili{Alor}
	and \ili{Pantar} which show only a slightly higher Papuan component.
	Towards the west, the central
	and western part of Flores (Bena
	and Rampasasa) and Sumba show a smaller Papuan component than found further \ili{east}.
	On Sumba, the Austronesian component is the biggest.
	On Flores, the admixture is more diverse.'' \citep[22]{Fricke2019}
\end{quote}

It is important to remember that genes
and languages are not inherently linked to each other.
For example, in urban Mexico there are hundreds of
thousands of people of genetic Uto-Aztecan or Mayan
ancestry who speak Indo-European Spanish
and not their heritage languages.
Languages move across genetic populations.
Both individuals and whole societies can shift
to speak languages unrelated to their DNA ancestry.
So the DNA evidence noted above is useful for confirming
that different gene pools have mixed in the
target region, and that there is a genetic boundary
that roughly coincides with the Blust line,
but is not otherwise very useful linguistically.
To highlight a few key points:

\begin{itemize}
	\item There is evidence of human habitation in our target region dating back to 42,000--35,000 BP.
				These humans had boats or other technology that allowed them
				to make multiple deep-water crossings between
				islands that were tens of kilometres distant,
				and to catch large numbers of open ocean (pelagic) fish species.
				We do not know what languages these early inhabitants spoke.
				Given the time depth
				and the technology and the yearly \ili{east}
				and west monsoons, there is no reason to assume there was only one journey by one group at one time in one direction.
	\item There were groups speaking a variety of Papuan languages scattered across our target region
				at l\ili{east} as far west as eastern \ili{Sumbawa},
				prior to the arrival of the Austronesian-speaking peoples in the region around 4,000--3,500 BP.
	\item Also prior to the arrival of the Austronesian-speaking peoples,
				several important food crops were dispersing \it{westward} out of the New Guinea region.
				These included a species of banana,
				sugarcane, taro, breadfruit and sago.
				The spread of terms such as {\#}muku ‘banana’ across unrelated language families raises serious questions
				for the region about subgrouping arguments that are based on the scattered distribution of certain lexical items.
	\item When the Austronesian-speaking peoples arrived in the region,
				again there is no reason to assume there was only one journey by one group at one time in one direction.
				They had the technology to go in any direction at any time,
				seas and weather permitting.
				The current limited archaeological evidence has a 300 -- 500 year time difference between evidence for
				Austronesian-speaking peoples in the Timor region (earlier),
				and in North Maluku (later),
				which is quite different from the scenario proposed in \citet[98]{Blust1993,Blust2008}.
				\ili{As} we show in this study,
				the complexity and discontinuities in the linguistic data
				in the region suggests either very rapid incursion(s),
				or the possibility of multiple incursions,
				possibly at multiple times,
				possibly from several directions,
				from different starting points,
				and different regions where they got off their boats.
	\item We infer that the persistence over three or four millennia of Papuan languages in the region to
				the present day, along with structural, typological
				and lexical accommodation in both directions between Austronesian
				and Papuan languages, suggests a contact scenario of bilingualism
				and intermarriage over centuries that also resulted in some language shift,
				rather than necessarily a dominance model of the “farming/language dispersal” hypothesis.
				Warfare seems to have been mostly localised over land,
				women, and power, rather than large-scale territory-grabbing efforts.\footnote{
						See also \citet{Schapper2020a}
						that large-scale warfare in eastern Indonesia
						and Timor-Leste seems to post-date European contact.} 
	\item The human gene pool of our target region is “very complex
				and [{\ldots}] involved a great deal of admixture”,
				with a genetic transition near the Wallace line
				(\citealp{Oppenheimer2003}, \citealp[22]{Fricke2019}).
				While this is important for understanding the region,
				it does not tell us about the languages or their prehistory.

\end{itemize}
